\section{System overview and design principles}
\label{sec:system}

QDCVR (v2.3) is a self-hosted knowledge-base management system that turns a
folder of PDFs, Office documents, and scans into knowledge bases that agents
can be trusted to answer from. It is deployed as four cooperating services
behind one configuration file: a Python service that schedules parsing and
owns the OCR engine, a TypeScript web application that serves the interface
and acts as the server-side proxy, an MCP server that exposes the platform's
capabilities as tools \cite{mcp}, and storage engines holding documents,
vectors, graph edges and task state.

\subsection{Requirements}
\label{sec:req}

The requirements below were derived from the deployment rather than
postulated, and each names the failure it prevents.

\begin{description}
\item[R1 --- Parse heterogeneous enterprise documents.] Process regulations,
failure reports, scanned inspection records and journal articles arrive as
PDFs, Office documents and images. The system must convert them to text it
can govern, including scanned pages, tables and formulas, and it must make a
parse failure visible rather than storing an empty or garbled document. The
100-paper acquisition left 21 papers that never became retrievable, and that
gap is reported as a first-class result with its cause left unattributed
(Section~\ref{sec:e1}).
\item[R2 --- Organize by content, not by folder.] A document's home must be
decided by what it says: library membership, tags, descriptions and graph
edges are derived from the content.
\item[R3 --- Manage the collection over its lifetime.] Libraries and
documents are created, moved, renamed, merged, split, deleted and re-indexed.
Every such operation must preserve agreement between the document on disk,
the tree registry and the library metadata, verifiably after the fact.
\item[R4 --- Retrieve with several independent channels.] Lexical, dense,
catalog and graph channels fail differently, so the system must generate
candidates with each, under explicit budgets, and expose which channel
proposed what.
\item[R5 --- Do not promote a similarity score into a citation.] A candidate
becomes an answer only after its fetched content has been judged to contain
direct evidence for the question. When no candidate survives, the system
returns an explicit not-found report instead of the closest text.
\item[R6 --- Be operable and engine-agnostic.] The platform runs on a single
local workstation (RTX 5060, 8\,GB, Windows), degrades rather than crashes
when a model or a subprocess fails, exposes long-running work as pollable
tasks, and drives whichever agent harness the operator has configured.
\end{description}

\subsection{Design principles}
\label{sec:principles}

Six principles shape the architecture, each enforced in the implementation
rather than left to prompt discipline.

\begin{description}
\item[P1 --- Content-routed organization.] Documents are organized by what
they say, not by where they sat: placement, tagging, description and graph
edges are derived from the parsed content, and retrieval starts from a
library catalog rather than a directory path (\S\ref{sec:organization}).
\item[P2 --- Content outranks similarity.] No similarity score promotes a
document into an answer. A judging model must read the candidate's fetched
content and vote, and the platform fetches that content itself so an agent
cannot argue for a document it has not read (\S\ref{sec:gate}).
\item[P3 --- Every module terminates at a contract.] Ingestion ends in a
checked post-condition, retrieval in a budgeted return value, verification in
a rubric score, and management in an audit.
\item[P4 --- Fail closed.] Missing text, an engine error, a timeout or an
unparseable verdict rejects the candidate rather than passing it, and a
document that cannot be retrieved after indexing is a failed ingest rather
than a warning.
\item[P5 --- Leave an audit trail.] Every candidate, verdict, cut decision,
tool call and lifecycle operation is recorded, so a claim about the system
can be re-examined from evidence instead of re-argued from memory.
\item[P6 --- Engine-agnostic agency.] Answer agents are pluggable harnesses
(14 adapters, including \texttt{omp}, Claude and ACP-compatible CLIs); the
platform constrains them through tool whitelists, budgets and prompts, never
by coupling to a particular vendor.
\end{description}

\subsection{Architecture and subsystems}
\label{sec:architecture}

Fig.~\ref{fig:arch} shows the four layers. The \emph{client and agent layer}
contains the web interface and the chat endpoint that adapts external agent
harnesses; requests are annotated with a lane policy before they leave this
layer. The \emph{tool layer} is an MCP server that exposes the platform's
capabilities as 95 tools in twelve groups (knowledge-base catalog, document
read/write, file system, parsing and task status, tags, search, the
verification gate, vector index and repair, knowledge graph, experience,
persona, and platform lifecycle and health);
the lane policy is a server-side whitelist, so a lane physically cannot call
a tool outside its set. The \emph{domain service layer} holds the tree and
catalog services, the parsing scheduler that manages the OCR engine as a
supervised subprocess, the vector and two-stage candidate services, the
experience service, the graph service and the resident judging worker. The
\emph{storage layer} holds the document tree with its registries, the vector
collections, the graph database and the task and session databases.

Three properties of this arrangement matter later. First,
the retrieval lanes of Section~\ref{sec:lanes} are not separate programs:
they are policies over one tool layer, which is why a new capability is
available to every lane at once and why lane behaviour is comparable across
harnesses. Second, long-running work (parsing, indexing, graph construction,
distillation) is exposed as a pollable task rather than a blocking call, so a
slow operation cannot stall an agent's turn. Third, the judging service is a
resident worker for engineering reasons, not architectural elegance: loading
the verifier per request cost 10--12\,s against 1--2\,s for a warm verdict,
which is why the worker, its file lock and its verdict cache exist
(\S\ref{sec:impl}).

\begin{figure}[t]
\centering
\includegraphics[width=\textwidth,height=0.86\textheight,keepaspectratio]{fig1_architecture}
\caption{QDCVR architecture. Documents enter through the acquisition
pipeline (left) and are parsed, quality-gated, split, described and indexed
into four stores. At run time (right), clients and agent harnesses reach the
platform through the chat endpoint, which attaches a lane policy enforced
against the MCP tool layer (95 tools); all retrieval lanes terminate at the
content-verification gate. Organization and life-cycle management (bottom)
keep the tree registry, the document files and the library metadata in
agreement.}
\label{fig:arch}
\end{figure}

Table~\ref{tab:modules} maps each function of the system to the components
that implement it, the design decision that shapes it, and the experiment
that provides evidence for it. The table is the paper's reading guide:
Sections~\ref{sec:acquisition}--\ref{sec:implementation} describe these rows
in order, and Section~\ref{sec:eval} evaluates them.

\begin{table}[!htbp]
\centering
\footnotesize
\setlength{\tabcolsep}{3.5pt}
\caption{System functions, their implementing components, the decision that
shapes each one, and the evidence that covers it. Evidence identifiers
E1--E8 refer to the evaluation subsections of Section~\ref{sec:eval}.}
\label{tab:modules}
\setlength{\tabcolsep}{5pt}
\resizebox{\textwidth}{!}{%
\begin{tabular}{p{2.5cm}p{3.9cm}p{4.5cm}p{1.8cm}}
\toprule
Function & Components and representative tools & Key design decision & Evidence \\
\midrule
Acquisition and parsing & parsing scheduler, OCR subprocess, duplicate
detector, split gate & quality-gated: a document that cannot be parsed is
rejected visibly, and a document that cannot be retrieved afterwards is a
failed ingest & E1 \\
Knowledge organization & content router, catalog, tag registry, description
generator & content decides placement; the catalog is the first retrieval
instrument & E1, E2 \\
Life cycle and integrity & tree registry, library metadata, re-index tasks,
three-source audit & re-indexing is explicit; every operation is audited
afterwards & E2 \\
Document graph & graph projection of documents, libraries and tags;
related-document queries & the graph is a navigation index, never an evidence
source & E2 \\
Catalog-first navigation & library and document listings, progressive
disclosure & metadata before bodies, so a candidate set is chosen without
loading content & E4, E5 \\
Content-verified retrieval & vector, lexical, catalog and graph channels;
lane policies; retrieval gate & verification is a promotion condition, not a
ranking hint; verdicts fail closed & E4, E6 \\
Operational knowledge & experience store, credibility tiers, distillation
runs & a lesson is a claim about practice, so it is drafted, gated and
capped until it is vetted & E3 \\
Agent interface and operations & MCP tool layer, chat endpoint, lane
whitelists, task registry & every capability is a tool with a schema; long
work is pollable; engines are interchangeable & E5, E7 \\
Persona and platform lifecycle & persona distillation and persona Q\&A tools;
version and update tools & shipped and reachable through the same tool
surface, but outside this paper's evaluation, so no claim is made here & -- \\
\bottomrule
\end{tabular}
}
\end{table}

\subsection{Data model and storage layout}
\label{sec:datamodel}

The platform keeps its data model readable: documents are
Markdown files, and the two registries are JSON and YAML files that a human
can open. Because they are files rather than a database, a document's address is a path and its identity a path plus a content hash, which is what makes
the three-source audit of \S\ref{sec:org-registries} possible. The vector
store is a separate engine that can disagree with the registries; the system
treats that disagreement as a detectable, repairable condition rather than as
an impossibility. Table~\ref{tab:datamodel} lists the stores and fields the
paper refers to.



\subsection{Object life cycle}
\label{sec:lifecycle}

A document moves through five states. It is \emph{raw} when it arrives, with
a content hash recorded. It is \emph{parsed} when text has been extracted and
has passed the parse-quality checks. It is \emph{stored} when the body, its
description and its tags have been written to disk and registered in both
registries. It is \emph{indexed} when the vector store and the retrieval
indexes contain it, which requires an explicit indexing call. It is
\emph{verified} when a retrieval probe returns the document for a query
derived from its own content. Only the last state counts as a successful
ingest; a document that stalls at \emph{stored} is an incomplete operation,
not a warning.

Libraries have a simpler life cycle: a library is \emph{active} while it is
addressable and searchable, and \emph{merged} or \emph{split} when an
operator consolidates or divides it. Merge and split are orchestrated
operator procedures over the document, catalog and index tools rather than
single tool calls, and neither was exercised under controlled conditions
(Section~\ref{sec:e2}); we describe the transition without claiming it.
Library metadata travels with the
library when it moves, including its distillation configuration, which is why
the configuration is stored per library rather than in the deployment file.
Platform-level versioning is explicit: the platform reports its version
(2.3.0 in this deployment, the value carried by the repository's
\protect\path{VERSION} file and returned by
\protect\path{ragctl version --json}) and can update itself from its
repository through
the same command-line interface that performs the other operations. The
update path is the one platform-lifecycle probe that failed in E2: its check
timed out after 30\,s, so version reporting is evidenced and self-update is
not.

Re-indexing is where the object life cycle meets retrieval quality. It is
explicit, asynchronous and observable: it returns a task identifier, it can
be forced, and its post-condition is a segment count that matches the
document content. Nothing in the write path re-indexes as a side effect.
This trades convenience for the property that no document can silently be
missing from the vector store: the cost is that a writer who forgets the
call leaves a stale index, and the mitigation is the audit of
\S\ref{sec:org-integrity}, which detects exactly that condition.
